\begin{tresc}
Na rysunku poniżej widocznych jest 9 kółek: trzy białe, trzy czarne i trzy w kropki. Wykonując ruchy polegające na zamianie miejscami dwóch kółek chcemy doprowadzić do sytuacji, gdy w każdym wierszu i w każdej kolumnie będzie jedno kółko białe, jedno kółko czarne i jedno w kropki. Ile co najmniej ruchów musimy wykonać?
\rysunek{rysunek-24-stem-0.png}{Układ dziewięciu kółek}
\end{tresc}
\begin{wariant}{A} 1 \end{wariant}
\begin{wariant}{B} 2 \end{wariant}
\begin{wariant}{C} 3 \end{wariant}
\begin{wariant}{D} 4 \end{wariant}
\begin{wariant}{E} 5 \end{wariant}
